%! Trigonometry and Polar Coordinates — Unit 7, Lesson 7.6 — Exit Ticket (Answer Key)
\documentclass[10pt]{article}
\usepackage{precalculus-article}
\usepackage{precalculus-key}   % pulls in -boxes; do NOT also load -boxes
\newcommand{\TallMath}[1]{$\displaystyle #1\rule[-1.4em]{0pt}{3.2em}$}

\begin{document}

\pageheader{Unit 7, Lesson 7.6}{Exit Ticket --- Answer Key}
\namedateperiod

\vspace{0.12in}

\begin{notesbox}{Exit Ticket --- Key}

\textbf{1.}\quad Convert the polar coordinates $(4,\, \pi/3)$ to rectangular coordinates.

\vspace{6pt}
$x = 4\cos\!\left(\dfrac{\pi}{3}\right) = $ \ans{$4 \cdot \dfrac{1}{2} = 2$}

\vspace{8pt}
$y = 4\sin\!\left(\dfrac{\pi}{3}\right) = $ \ans{$4 \cdot \dfrac{\sqrt{3}}{2} = 2\sqrt{3}$}

\vspace{6pt}
Rectangular form: \ans{$(2,\; 2\sqrt{3})$}

\vspace{0.18in}
\textbf{2.}\quad Convert the rectangular coordinates $(-3,\, 3)$ to polar form
with $r > 0$ and $0 \le \theta < 2\pi$.

\vspace{6pt}
$r = \sqrt{(-3)^2 + 3^2} = $ \ans{$\sqrt{18} = 3\sqrt{2}$}

\vspace{8pt}
$\theta = $ \ans{$\arctan(3/(-3)) + \pi = \arctan(-1) + \pi = -\pi/4 + \pi = 3\pi/4$
\quad (point is in Q2, so add $\pi$)}

\vspace{6pt}
Polar form: \ans{$(3\sqrt{2},\; 3\pi/4)$}

\vspace{0.18in}
\textbf{3.}\quad Write the complex number $2 - 2i$ in polar form $(r\cos\theta) + i(r\sin\theta)$.

\vspace{6pt}
$r = $ \ans{$\sqrt{4+4} = 2\sqrt{2}$} \qquad
$\theta = $ \ans{$\arctan(-2/2) = -\pi/4$, or $7\pi/4$ in $[0,2\pi)$}

\vspace{6pt}
Polar form: \ans{$2\sqrt{2}\cos(7\pi/4) + 2\sqrt{2}\,i\sin(7\pi/4)$}

\end{notesbox}

\end{document}
